\section{Limits}
\label{sec:limits}

\paragraph{One corpus.}
LoCoMo is ten conversations. It is the standard benchmark for this
task~\cite{locomo} and its evidence annotations are what make the \textbf{all}
metric computable at all, but ten conversations is ten conversations, and the
structural facts in \S\ref{sec:diagnostic} --- particularly the
\rSpanSessions{}\% cross-session figure --- are properties of \emph{this} corpus.
A corpus whose multi-hop evidence clustered within sessions would very likely
behave differently, and the policies might help there.

\paragraph{One embedding.}
Every row uses \texttt{zen-embedding-0.6b}. The comparison between policies is
sound because they share an index, but the absolute recall numbers are that
model's, and a stronger embedding would shift the whole table. What would
\emph{not} shift is the relationship the paper is about: \textbf{any} exceeding
\textbf{all} by a wide margin is a statement about evidence distribution, not
about embedding quality.

\paragraph{Similarity standing in for reasoning.}
Our \textbf{chain} policy follows IRCoT's shape~\cite{ircot} with one important
substitution: IRCoT interleaves a language model's reasoning between retrieval
steps, and we substitute vector similarity for that step. This makes the policy
thousands of times cheaper and it is a fair test of ``does searching again from
what you found help'', but it is \emph{not} a test of IRCoT. A model that reads
the first hop and writes a genuinely new query has access to exactly the entity
linkage \S\ref{sec:diagnostic} identifies as missing. We expect it to do better;
we did not measure it; and it costs a model call per hop, which is a different
economic proposition from \rPolChainMs{}\,ms.

\paragraph{Retrieval only.}
We measure whether the evidence was retrieved, never whether a model answered
correctly given it. End-to-end accuracy depends on the reader as much as the
retriever, and conflating the two is how retrieval failures get attributed to
models. The scope here is deliberate and narrow.
